\exercisesection{Complexes of A-Modules}

\begin{enumerate}
\def\labelenumi{\arabic{enumi})}
\tightlist
\item
  Assume the ring A is principal; let C be a complex of free A-modules.
\end{enumerate}

\begin{enumerate}
\def\labelenumi{\alph{enumi})}
\item
  Show that C is the direct sum of a family of complexes \((^{p}\mathbf{E})_{p\in \mathbb{Z}}\) satisfying: \(^{p}\mathbf{E}_{n} = 0\) for \(n\neq p\) , \(p - 1 ; Z_{n}(^{p}\mathbf{E}) = 0\) .
\item
  We further assume that \(C_n\) is finitely generated for all \(n \in \mathbf{Z}\) . Show that \(C\) is a direct sum of complexes having only either one nonzero component, or two consecutive nonzero components, these components being isomorphic to \(A\) .
\end{enumerate}

\begin{enumerate}
\def\labelenumi{\arabic{enumi})}
\setcounter{enumi}{1}
\item
  Let C and C' be two complexes of A-modules; suppose that the complexes C and B(C) (resp. C' and Z(C')) are projective (resp. injective). Show that every morphism of complexes from H(C) to H(C') is induced by a morphism from C to C'.
\item
  Let C be a complex of A-modules; one regards C as an A(ε)-module (X, p. 27). We denote by
\end{enumerate}

\[
i: \mathbf {A} \rightarrow \mathbf {A} (\varepsilon)
\]

the canonical injection.

Show that the following conditions are equivalent:

\begin{enumerate}
\def\labelenumi{(\roman{enumi})}
\item
  C is i-projective (X, p. 170, exercise 19).
\item
  It is i-injective (loc. cit.).
\item
  The complex C is homotopic to zero.
\item
  There exists a graded A-module M and an isomorphism of graded A(ε)-modules A(ε) ⊗A M → C.
\item
  There exists a graded A-module N and an isomorphism of graded A(ε)-modules C → Hom\_A(A(ε), N).
\end{enumerate}

\begin{enumerate}
\def\labelenumi{\arabic{enumi})}
\setcounter{enumi}{3}
\tightlist
\item
  \begin{enumerate}
  \def\labelenumii{\alph{enumii})}
  \tightlist
  \item
    Let C be a right-bounded A-complex, such that C and H(C) are projective. Show that C is split.
  \end{enumerate}
\end{enumerate}

\begin{enumerate}
\def\labelenumi{\alph{enumi})}
\setcounter{enumi}{1}
\tightlist
\item
  On the ring B = A(ε) (X, p. 27), we consider the complex C defined by C\_n = B for all n, and d(b) = εb for b ∈ B. The complex C is free with zero homology, but not split.
\end{enumerate}

\begin{enumerate}
\def\labelenumi{\arabic{enumi})}
\setcounter{enumi}{4}
\tightlist
\item
  Let C be a complex, and let s be a splitting of C, \(\delta: \mathbf{C} \to \mathbf{C}\) a graded homomorphism of degree -1, such that \((d + \delta) \circ (d + \delta) = 0\) . One denotes \(\mathbf{C}'\) the complex \((\mathbf{C}, d + \delta)\) . Assume that the A-linear map \(\theta\) from C to C defined by \(\theta = 1_C + \delta s\) is bijective.
\end{enumerate}

\begin{enumerate}
\def\labelenumi{\alph{enumi})}
\item
  In order for \(s\theta^{-1}\) to be a splitting of \(C'\) , it is necessary and sufficient that \(\delta(Z(C)) \subset \theta(B(C))\) .
\item
  Assume the conditions of a) are satisfied. Show that there exist direct sum decompositions: \(\mathbf{C} = \mathbf{B}(\mathbf{C}) \oplus \mathbf{Ker}(ds) = \mathbf{B}(\mathbf{C}') \oplus \mathbf{Ker}(ds)\) . If p (resp. \(p'\) ) is the projector with image \(\mathbf{Ker}(ds)\) and kernel \(\mathbf{B}(\mathbf{C})\) (resp. \(\mathbf{B}(\mathbf{C}')\) ), we have \(p' = p\theta^{-1}\) .
\item
  We further assume that the A-linear map \(\lambda\) from C to C defined by \(\lambda = 1_{C} + s\delta\) is bijective. Show that we have \(Z(C') = \lambda^{-1}(Z(C))\) and \(C = Z(C) \oplus \operatorname{Im}(sd) = Z(C') \oplus \operatorname{Im}(sd)\) . If \(q\) (resp. \(q'\) ) is the projector with image \(Z(C)\) (resp. \(Z(C')\) ) and kernel \(\operatorname{Im}(sd)\) , we have \(q' = \lambda^{-1} q\) .
\end{enumerate}

\begin{enumerate}
\def\labelenumi{\arabic{enumi})}
\setcounter{enumi}{5}
\tightlist
\item
  We keep the notation from Exercise 5; we further assume that A is commutative, that the complexes C and \(\mathbf{C}^{\prime}/\mathbf{B}(\mathbf{C}^{\prime})\) are flat and that there exists an ideal m of A satisfying \(\delta(\mathbf{C}) \subset \mathfrak{m}\mathbf{C}\) and \(\bigcap_{n} (\mathbf{B}(\mathbf{C}) + \mathfrak{m}^{n}\mathbf{C}) = \mathbf{B}(\mathbf{C})\) (This last condition is always satisfied if A is a local Noetherian ring with maximal ideal m and if \(C_{n}\) is a finitely generated A-module for every n, cf. AC, III, § 3, n° 3, prop. 6). Then show that \(s\theta^{-1}\) is a splitting of \(C^{\prime}\) .
\end{enumerate}

(The flatness hypotheses imply \(\operatorname{Im}(d + \delta) \cap \mathfrak{m}^n\mathbf{C} = (d + \delta)\, (\mathfrak{m}^n\,\mathbf{C})\) for all \(n\) ; if \(x \in \mathbf{Z}(\mathbf{C})\) , construct by induction elements \(y_n \in \mathbf{C}\) , \(z_n \in \mathbf{Z}(\mathbf{C}) \cap \mathfrak{m}^n\mathbf{C}\) such that \(\delta(x) = \theta d(y_n) + \delta(z_n)\) . Conclude using exercise 5, a).)

\begin{enumerate}
\def\labelenumi{\arabic{enumi})}
\setcounter{enumi}{6}
\tightlist
\item
  Let C and C' be two complexes, \(s_{1}(\text{resp. } s')\) a splitting of C (resp. C'), \(u : C' \to C\) a morphism of complexes. We define an A-endomorphism \(\sigma\) of Con(u), graded of degree -1, by setting
\end{enumerate}

\[
\sigma (y ^ {\prime}, x) = \left(- s ^ {\prime} (y ^ {\prime}), s (x) - s u s ^ {\prime} (y ^ {\prime})\right).
\]

Show that in order for \(\sigma\) to be a splitting of Con \((u)\) , it is necessary and sufficient that we have H(u) = 0.

\begin{enumerate}
\def\labelenumi{\arabic{enumi})}
\setcounter{enumi}{7}
\tightlist
\item
  Let \(0 \to M' \xrightarrow{\mu} M \xrightarrow{\nu} M'' \to 0\) an exact sequence of A-modules, which can be viewed as a sequence
\end{enumerate}

exact sequence of complexes (zero in degrees ≠ 0). Show that the morphism of complexes φ : Con(u) → M''\,(X, p. 39) is a homotopy equivalence if and only if the exact sequence 0 → M' → M → M'' → 0 is split. 9) Let u : C' → C be a morphism of A-complexes. Denote by v : C → Coker(u) the quotient map. Consider the graded A-homomorphisms

\[
\alpha : (\operatorname{Ker} u) (- 1) \rightarrow \operatorname{Con} (u) \quad \text { and } \quad \varphi : \operatorname{Con} (u) \rightarrow \operatorname{Coker} (u)
\]

defined by \(\alpha(x') = (x', 0)\) , \(\varphi(x', x) = v(x)\) for \(x' \in C'\) , \(x \bullet C\) . Show that \(\alpha, \varphi\) are morphisms of complexes and that they give rise to an exact sequence

\[
\dots \rightarrow \mathrm{H} _ {n - 1} (\text { Ker } u) \xrightarrow {\mathrm{H} (\alpha)} \mathrm{H} _ {n} (\text { Con } (u)) \xrightarrow {\mathrm{H} (\varphi)} \mathrm{H} _ {n} (\text { Coker } u) \xrightarrow {\delta} \mathrm{H} _ {n - 2} (\text { Ker } u) \rightarrow \dots
\]

where δ is the composite of the connecting homomorphisms relative to the two exact sequences

\[
0 \rightarrow \operatorname{Ker} u \rightarrow C ^ {\prime} \rightarrow \operatorname{Im} u \rightarrow 0 \quad \text { and } \quad 0 \rightarrow \operatorname{Im} u \rightarrow C \rightarrow \operatorname{Coker} u \rightarrow 0.
\]

\begin{enumerate}
\def\labelenumi{\arabic{enumi})}
\setcounter{enumi}{9}
\tightlist
\item
  Let \(u: \mathbf{C} \to \mathbf{C}'\) a morphism of A-complexes, \(\mathbf{P}\) a projective A-module, \(n\) be an integer \(\geqslant 0\) ,
\end{enumerate}

\[
h: \mathrm{P} \rightarrow \mathrm{H} _ {n} (\text { Con } (u))
\]

an A-homomorphism. Show that there exists a complex \(\tilde{C}\) , an exact sequence of complexes

\[
0 \to \mathbf {C} \xrightarrow {i} \tilde {\mathbf {C}} \to \mathbf {P} (- n) \to 0
\]

and a morphism \(\tilde{u}:\tilde{\mathbf{C}}\to \mathbf{C}'\) such that \(\tilde{u}\circ i = u\) and that the morphism \(\mathbf{C}(i):\mathrm{Con}(u)\to \mathrm{Con}(\tilde{u})\) induced by \(i\) satisfies the following properties:

α) The map \(\mathrm{H}_{p}(\mathrm{C}(i)) : \mathrm{H}_{p}(\mathrm{Con}(u)) \to \mathrm{H}_{p}(\mathrm{Con}(\tilde{u}))\) is bijective for \(p \neq n, n + 1\) ;

β) There exists an exact sequence

\[
0 \rightarrow \mathrm{H} _ {n + 1} (\operatorname{Con} (u)) \xrightarrow {\mathrm{H} _ {n + 1} (\mathbf {C} (i))} \mathrm{H} _ {n + 1} (\operatorname{Con} (\tilde {u})) \rightarrow \mathrm{P} \xrightarrow {h} \mathrm{H} _ {n} (\operatorname{Con} (u)) \xrightarrow {\mathrm{H} _ {n} (\mathbf {C} (i))} \mathrm{H} _ {n} (\operatorname{Con} (\tilde {u})) \rightarrow 0.
\]

\begin{enumerate}
\def\labelenumi{\arabic{enumi})}
\setcounter{enumi}{10}
\tightlist
\item
  We call a bicomplex of A-modules a triple (C, \(d'\) , \(d''\) ), where C is a graded A-module of type \(\mathbf{Z} \times \mathbf{Z}\) , \(d': \mathbf{C} \to \mathbf{C}\) a graded A-endomorphism of degree (-1, 0) and \(d'': \mathbf{C} \to \mathbf{C}\) a graded A-endomorphism of degree (0, -1), satisfying \(d' \circ d' = d'' \circ d'' = d' \circ d'' + d'' \circ d' = 0\) We will also use the ascending graduation defined by \(\mathbf{C}^{p,q} = \mathbf{C}_{-p,-q}\) for \(p, q \in \mathbf{Z}\) .
\end{enumerate}

We denote by \(Z'(C)\) , \(Z''(C)\) , \(B'(C)\) , \(B''(C)\) the sub-bicomplexes \(\text{Ker } (d')\) , \(\text{Ker } (d'')\) , \(\text{Im } (d')\) , \(\text{Im } (d'')\) , and \(H'(C)\) (resp. \(H''(C)\) ) the bicomplex \(Z'(C)/B'(C)\) (resp. \(Z''(C)/B''(C)\) ).

\begin{enumerate}
\def\labelenumi{\alph{enumi})}
\tightlist
\item
  We set \(d = d' + d''\) . Show that when C is equipped with the total graduation associated with the given bigraduation (defined by \(\mathbf{C}_n = \bigoplus_{p+q=n} \mathbf{C}_{p,q}\) , cf. II, p. 164), the pair (C, d) is a complex. This is called
\end{enumerate}

the complex associated with the bicomplex (C, \(d'\) , \(d''\) ); its homology is denoted H(C).

\begin{enumerate}
\def\labelenumi{\alph{enumi})}
\setcounter{enumi}{1}
\item
  Let (C, \(d'\) , \(d''\) ) and (K, \(\delta'\) , \(\delta''\) ) be two bicomplexes. A morphism from (C, \(d'\) , \(d''\) ) to (K, \(\delta'\) , \(\delta''\) ) is a graded A-homomorphism \(u\) of degree (0, 0) from C to K satisfying: \(\delta'\circ u = u \circ d'\) and \(\delta''\circ u = u \circ d''\) . Show that \(u\) induces a morphism of the associated complexes.
\item
  Let \(u, v\) two morphisms of bicomplexes from C to K. One calls a homotopy connecting \(u\) to \(v\) a pair \((s', s'')\) of A-homomorphisms from C to K, graded of degrees (1, 0) and (0, 1) respectively, such that
\end{enumerate}

\[
g - f = d ^ {\prime} \circ s ^ {\prime} + s ^ {\prime} \circ d ^ {\prime} + d ^ {\prime \prime} \circ s ^ {\prime \prime} + s ^ {\prime \prime} \circ d ^ {\prime \prime}; \quad s ^ {\prime} \circ d ^ {\prime \prime} + d ^ {\prime \prime} \circ s ^ {\prime} = 0; \quad s ^ {\prime \prime} \circ d ^ {\prime} + d ^ {\prime} \circ s ^ {\prime \prime} = 0.
\]

Show that \(s' + s''\) is a homotopy connecting the associated complex morphisms induced by \(u\) and \(v\) . 12) Let \((\mathbf{C}, d', d'')\) a bicomplex, such that one has \(\mathrm{H}_{p,-p}''(\mathbf{C}) = 0\) for all \(p \in \mathbf{Z}\) , \(\mathrm{H}_{-q,q+1}'(\mathbf{C}) = 0\) for \(q \geqslant 0\) , \(\mathrm{H}_{r,-r-1}'(\mathbf{C}) = 0\) for \(r \geqslant 0\) , and that there exist positive integers \(a\) and \(b\) such that \(\mathbf{C}_{a,-a} = \mathbf{C}_{-b,b} = 0\) . Prove that the A-module \(\mathrm{H}_{0,0}'(\mathbf{C})\) is zero.

\begin{enumerate}
\def\labelenumi{\arabic{enumi})}
\setcounter{enumi}{12}
\tightlist
\item
  Let I be a finite set; we denote by \((e_i)_{i\in I}\) the canonical basis of the Z-module \(\mathbf{Z}^{(I)}\) . We call an I-complex (resp. I-precomplex) of A-modules a graded A-module C of type \(\mathbf{Z}^{I}\) , equipped with a family of A-endomorphisms \((d_i)_{i\in I}\) , such that \(d_i\) be graded of degree \((-e_i)\) satisfying:
\end{enumerate}

\[
d _ {i} \circ d _ {i} = 0 \quad \text { for every } \quad i \in I;
\]

\[
d _ {i} \circ d _ {j} + d _ {j} \circ d _ {i} = 0 (\text { resp. } d _ {i} \circ d _ {j} = d _ {j} \circ d _ {i}) \quad \text { for } \quad i, j \in I.
\]

If Card (I) = 1, an I-complex is a complex; if Card (I) = 2, an I-complex is a bicomplex (exercise 11). If \((\mathbf{C}, (d_i))\) and \((\mathbf{C}', (d_i'))\) are two I-complexes (resp. I-precomplexes), a morphism of \((\mathbf{C}, (d_i))\) to \((\mathbf{C}', (d_i'))\) is a graded homomorphism of degree zero \(u: \mathbf{C} \to \mathbf{C}'\) such that \(d_i' \circ u = u \circ d_i\) for all \(i \in I\) .

\begin{enumerate}
\def\labelenumi{\alph{enumi})}
\tightlist
\item
  Let (C, \((d_i)\) ) an I-complex. If one equips C with the total grading (of type Z) deduced from its grading, show that (C, \(\sum_{i\in I}d_i\) ) is a complex of A-modules, called the complex associated with the I-complex (C, \((d_i)\) ).
\end{enumerate}

Show that a morphism of I-complexes induces a morphism of the associated complexes.

\begin{enumerate}
\def\labelenumi{\alph{enumi})}
\setcounter{enumi}{1}
\tightlist
\item
  Let C be an I-precomplex, and let R be a total order relation on I, denoted ≤. For \(k = (k_i)_{i \in I}, x \in C_k\) and \(i \in I\) , we set
\end{enumerate}

\(\delta_{i}(x) = (-1)^{\sum_{j < i} k_{j}} d_{i}(x)\) . Show that \((C, (\delta_{i}))\) is an I-complex, denoted \(C_{R}\) .

\begin{enumerate}
\def\labelenumi{\alph{enumi})}
\setcounter{enumi}{2}
\item
  Let \(\mathbf{R}_{\alpha}, \mathbf{R}_{\beta}\) two total order relations on I, denoted \(\leqslant_{\alpha}\) and \(\leqslant_{\beta}\) . Let \(u_{\beta \alpha}\) the endomorphism of C defined by \(u(x) = (-1)^{\varepsilon(k)} x\) for \(x \in \mathbf{C}_k\) , with \(\varepsilon(k) = \sum_{\substack{i < _{\alpha} j \\ j < _{\beta} i}} k_i k_j\) . Show that \(u_{\beta \alpha}\) is a morphism of I-complexes of \(\mathbf{C}_{\mathbf{R}_{\alpha}}\) to \(\mathbf{C}_{\mathbf{R}_{\beta}}\) .
\item
  Let S be the set of total order relations on I; we endow S with the trivial preorder (that is, the preorder such that \(R_{\alpha} \leqslant R_{\beta}\) whatever \(R_{\alpha}, R_{\beta} \in S\) ). Show that the system \((\mathbf{C}_{R\alpha}, u_{\alpha\beta})\) is a projective system of I-complexes, relative to S; its limit C is an I-complex, called the I-complex associated with the I-precomplex C. Show that every morphism of I-precomplexes induces a morphism of the associated I-complexes.
\end{enumerate}

A spectral sequence of A-modules is defined as a sequence of A-complexes \((\mathrm{E}_r, d_r)_{r \geqslant 2}\) and of mappings \(\alpha_r: \mathrm{H}(\mathrm{E}_r) \to \mathrm{E}_{r+1}\) such that:

\begin{enumerate}
\def\labelenumi{(\roman{enumi})}
\tightlist
\item
  The grading of the complex \(E_r\) is the total grading associated with a bigrading
\end{enumerate}

\[
\mathrm{E} _ {r} = \bigoplus_ {p, q \in \mathbf {Z}} \mathrm{E} _ {r} ^ {p, q} (r \geqslant 2);
\]

we have \(d_r(E_r^{pq}) \subset E_r^{p + r,q - r + 1}\) .

\begin{enumerate}
\def\labelenumi{(\roman{enumi})}
\setcounter{enumi}{1}
\tightlist
\item
  The maps \(\alpha_{r}\) are isomorphisms of bigraded A-modules.
\end{enumerate}

We will also use the descending bigrading of \(E_{r}\) : we denote \(E_{pq}^{r} = E_{r}^{-p, -q}\) .

\begin{enumerate}
\def\labelenumi{\alph{enumi})}
\tightlist
\item
  Here we denote by \(Z_{r+1}(E_r)\) (resp. \(B_{r+1}(E_r)\) ) the bigraded A-module Ker \(d_r\) (resp. Im \(d_r\) ). Let
\end{enumerate}

\[
p _ {r}: \mathrm{E} _ {r} \rightarrow \mathrm{E} _ {r} / \mathrm{B} _ {r + 1} (\mathrm{E} _ {r})
\]

the quotient map and \(i_{r}: \mathbf{E}_{r+1} \to \mathbf{E}_{r}/\mathbf{B}_{r+1}(\mathbf{E}_{r})\) the composite map of \(\alpha_{r}^{-1}\) and the natural injection. For every \(r \geqslant 2\) , we define by induction on \(k\) the bigraded submodules \(\mathbf{Z}_{r+k}(\mathbf{E}_{r})\) and \(\mathbf{B}_{r+k}(\mathbf{E}_{r})\) of \(\mathbf{E}_{r}\) by setting:

\[
\mathbf {Z} _ {r + k} (\mathrm{E} _ {r}) = p _ {r} ^ {- 1} \left(i _ {r} \left(\mathbf {Z} _ {r + k} (\mathrm{E} _ {r + 1})\right)\right); \quad \mathbf {B} _ {r + k} (\mathrm{E} _ {r}) = p _ {r} ^ {- 1} \left(i _ {r} \left(\mathbf {B} _ {r + k} (\mathrm{E} _ {r + 1})\right)\right).
\]

Prove that we have the inclusions:

\[
0 \subset \mathrm{B} _ {r + 1} (\mathrm{E} _ {r}) \subset \mathrm{B} _ {r + 2} (\mathrm{E} _ {r}) \subset \dots \subset \mathrm{Z} _ {r + 2} (\mathrm{E} _ {r}) \subset \mathrm{Z} _ {r + 1} (\mathrm{E} _ {r}) \subset \mathrm{E} _ {r}
\]

and that \(\mathbf{E}_k\) is isomorphic to \(\mathbf{Z}_k(\mathbf{E}_r) / \mathbf{B}_k(\mathbf{E}_r)\) for \(k\geqslant r + 1\)

We set \(Z_{\infty}(E_2) = \bigcap_{r \geqslant 3} Z_r(E_2), B_{\infty}(E_2) = \bigcup_{r \geqslant 3} B_r(E_2)\) and \(E_{\infty} = Z_{\infty}(E_2)/B_{\infty}(E_2)\) . We say that the spectral sequence is regular if the decreasing sequence \((Z_r(E_2^{pq}))_{r \geqslant 2}\) is stationary for every pair \((p, q)\) .

\begin{enumerate}
\def\labelenumi{\alph{enumi})}
\setcounter{enumi}{1}
\item
  Let \(E = (E_r, d_r, \alpha_r)_{r \geqslant 2}\) and \(E' = (E_r', d_r', \alpha_r')_{r \geqslant 2}\) two spectral sequences of A-modules; a morphism from E to \(E'\) is a sequence of morphisms of A-complexes \(u_r : E_r \to E_r'\) , compatible with the bigradings, such that \(\alpha_r' \circ H(u_r) = u_{r+1} \circ \alpha_r\) for all \(r \geqslant 2\) Show that for every \(r \geqslant 2\) and for every \(k \geqslant r + 1\) , \(u_r\) induces a homomorphism of \(Z_k(E_r)\) (resp. \(B_k(E_r)\) in \(Z_k(E_r')\) (resp. \(B_k(E_r')\) ). In particular, \(u_2\) induces homomorphisms \(Z_\infty(u_2) : Z_\infty(E_2) \to Z_\infty(E_2')\) , \(B_\infty(u_2) : B_\infty(E_2) \to B_\infty(E_2')\) and \(u_\infty : E_\infty \to E_\infty'\) ; all these homomorphisms are compatible with the bigradings.
\item
  We assume that \(u_2\) is an isomorphism. Show that then \(u_r\) is an isomorphism for all \(r \geqslant 2\) as well as \(\mathbf{B}_{\infty}(u_2)\) . If moreover the sequences E and E' are regular, \(\mathbf{Z}_{\infty}(u_2)\) and \(u_{\infty}\) are isomorphisms.
\item
  Let \(G = \bigoplus_{n \in Z} G^n\) a graded A-module, equipped with a decreasing sequence \((F^p G)_{p \in \mathbb{Z}}\) of graded sub-A-modules
\end{enumerate}

such that \(\bigcap_{p} F^{p} G = 0\) and \(\bigcup_{p} F^{p} G = G\) . We say that the spectral sequence E approaches \((G, (F^{p} G)_{p \in \mathbb{Z}})\)

(or simply G) if there exist isomorphisms \(\beta^{pq}: E_{\infty}^{pq} \to (F^p G)^{p+q}/(F^{p+1} G)^{p+q}\) for all

\[
(p, q) \in \mathbf {Z} \times \mathbf {Z}.
\]

One says that the spectral sequence E converges to G if it approaches G and if moreover:

\begin{enumerate}
\def\labelenumi{(\roman{enumi})}
\item
  The spectral sequence E is regular;
\item
  For every integer n, there exists an integer p such that \((\mathbf{F}^{p}\mathbf{G})^{n}=0\) .
\end{enumerate}

Suppose that E (resp. E′) converges to G (resp. G′). Let \(u: \mathrm{E} \to \mathrm{E}'\) a morphism of spectral sequences, \(v: \mathrm{G} \to \mathrm{G}'\) a graded homomorphism of degree 0, such that \(v(\mathrm{F}^p \mathrm{G}) \subset \mathrm{F}^p \mathrm{G}'\) for all \(p\) and that the induced homomorphisms \(\bar{v}^{p,n}: (\mathrm{F}^p \mathrm{G})^n / (\mathrm{F}^{p+1} \mathrm{G})^n \to (\mathrm{F}^p \mathrm{G}')^n / (\mathrm{F}^{p+1} \mathrm{G}')^n\) satisfy \(\bar{v}^{p,p+q} \circ \beta^{p,q} = \beta'^{p,q} \circ u_{\infty}^{p,q}\) for every pair \((p, q)\) . Show that if \(u_2^{pq}\) is an isomorphism for every pair \((p, q)\) , then \(v\) is an isomorphism.

\begin{enumerate}
\def\labelenumi{\arabic{enumi})}
\setcounter{enumi}{14}
\tightlist
\item
  Let E be a spectral sequence of A-modules, which approaches a graded A-module G (exercise 14).
\end{enumerate}

\begin{enumerate}
\def\labelenumi{\alph{enumi})}
\tightlist
\item
  Assume \(E_{2}^{p,q}=0\) for p\textless0. Show that the maps \(\alpha_{r}\) induce for every q injective A-homomorphisms:
\end{enumerate}

\[
\mathrm{E} _ {\infty} ^ {0, q} \rightarrow \dots \rightarrow \mathrm{E} _ {3} ^ {0, q} \rightarrow \mathrm{E} _ {2} ^ {0, q}
\]

whence, by composition with \((\beta^{0q})^{-1}\) and with the passage to the quotient map, an A-homomorphism \(\gamma^q: G^q \to E_2^{0,q}\) .

\begin{enumerate}
\def\labelenumi{\alph{enumi})}
\setcounter{enumi}{1}
\tightlist
\item
  Assume \(E_{2}^{pq} = 0\) for q \textless{} 0. Show that the maps \(\alpha_{r}^{-1}\) induce surjective homomorphisms:
\end{enumerate}

\[
\mathrm{E} _ {2} ^ {p, 0} \rightarrow \mathrm{E} _ {3} ^ {p, 0} \rightarrow \dots \rightarrow \mathrm{E} _ {\infty} ^ {p, 0}
\]

whence, by composition with \(\beta^{p,0}\) , an A-homomorphism \(\delta^p:E_2^{p,0}\to G^p\) c) We assume that the hypotheses of a) and b) are satisfied, that is, \(E_{2}^{pq} = 0\) if p \textless{} 0 or q \textless{} 0. Show that the sequence

\[
0 \rightarrow \mathrm{E} _ {2} ^ {1, 0} \xrightarrow {\delta^ {1}} \mathrm{G} ^ {1} \xrightarrow {\gamma^ {1}} \mathrm{E} _ {2} ^ {0, 1} \xrightarrow {d _ {2}} \mathrm{E} _ {2} ^ {2, 0} \xrightarrow {\delta^ {2}} \mathrm{G} ^ {2}
\]

is exact.

\begin{enumerate}
\def\labelenumi{\alph{enumi})}
\setcounter{enumi}{3}
\item
  Assume \(\mathbf{E}_2^{pq} = 0\) if \(p > 0\) or \(q > 0\) ; define an exact sequence analogous to that of \(c\) ).
\item
  One assumes that there exist integers \(r, d\) , with \(d \geqslant r \geqslant 2\) , such that \(\mathrm{E}_{r}^{pq} = 0\) if \(p \neq 0\) , \(d\) . Define an exact sequence:
\end{enumerate}

\[
\dots \rightarrow \mathrm{E} _ {r} ^ {d, n - d} \rightarrow \mathrm{G} ^ {n} \rightarrow \mathrm{E} _ {r} ^ {0, n} \rightarrow \mathrm{E} _ {r} ^ {d, n + 1 - d} \rightarrow \mathrm{G} ^ {n + 1} \rightarrow \dots .
\]

If likewise \(E_{r}^{pq} = 0\) for \(q \neq 0\) , d, show that there exists an exact sequence:

\[
\dots \rightarrow \mathrm{E} _ {r} ^ {n, 0} \rightarrow \mathrm{G} ^ {n} \rightarrow \mathrm{E} _ {r} ^ {n - d, d} \rightarrow \mathrm{E} _ {r} ^ {n + 1, 0} \rightarrow \dots .
\]

\begin{enumerate}
\def\labelenumi{\alph{enumi})}
\setcounter{enumi}{5}
\tightlist
\item
  Assume \(\mathbf{E}_2^{pq} = 0\) if \(q \neq 0\) (resp. \(p \neq 0\) ). Show that the homomorphism
\end{enumerate}

\[
\delta^ {n}: \mathrm{E} _ {2} ^ {n, 0} \rightarrow \mathrm{G} ^ {n} (\text { resp. } \gamma^ {n}: \mathrm{G} ^ {n} \rightarrow \mathrm{E} _ {2} ^ {0, n})
\]

is bijective for every n.

¶ 16) Let C be an A-complex, \((F^{p}C)_{p\in\mathbf{Z}}\) a decreasing sequence of subcomplexes of C.

\begin{enumerate}
\def\labelenumi{\alph{enumi})}
\tightlist
\item
  For \(p,q,r\in\mathbf{Z}\) , let \(M_{r}^{pq}\) the set of elements x of \((F^{p}C)^{p+q}\) such that \(dx\in F^{p+r}C\) ; for \(r\geqslant0\) , we set \(E_{r}^{pq}=M_{r}^{pq}/(d(M_{r-1}^{p-r+1,q+r-2})+M_{r}^{p+1,q-1})\) and we denote by \(d_{r}:E_{r}^{pq}\to E_{r}^{p+r,q-r+1}\) the homomorphism deduced from \(d\) by passing to quotients. Show that there exists for \(r\geqslant0\) an isomorphism \(\alpha_{r}:H(E_{r})\to E_{r+1}\) , so that the sequence \((E_{r},d_{r})_{r\geqslant2}\) , equipped with the maps \(\alpha_{r}\) , is a spectral sequence. We have
\end{enumerate}

\[
\mathrm{E} _ {0} ^ {p q} = (\mathrm{F} ^ {p} \mathrm{C}) ^ {p + q} / (\mathrm{F} ^ {p + 1} \mathrm{C}) ^ {p + q}.
\]

\begin{enumerate}
\def\labelenumi{\alph{enumi})}
\setcounter{enumi}{1}
\item
  We assume that \(\bigcup_{p} \mathbf{F}^{p} \mathbf{C} = \mathbf{C}\) and that for every \(n\) , there exists an integer \(p\) such that \((\mathbf{F}^{p} \mathbf{C})^{n} = 0\) . Show that the spectral sequence defined in \(a)\) converges (exercise 14) to the graded A-module \(\mathrm{H}(\mathbf{C})\) , equipped with the sequence of graded submodules \(\mathbf{F}^{p} \mathbf{H}(\mathbf{C}) = \operatorname{Im} \mathbf{H}(i_{p})\) where \(i_{p}: \mathbf{F}^{p} \mathbf{C} \to \mathbf{C}\) is the canonical injection.
\item
  If \(\mathbf{F}^0\bar{\mathbf{C}} = \mathbf{C}\) the A-module \(\mathbf{E}_r^{pq}\) is zero for \(p < 0\) ; if \((\mathbf{F}^p\mathbf{C})^n = 0\) for \(p > n\) , we have \(\mathbf{E}_r^{pq} = 0\) for \(q < 0\) . d) Let \(\mathbf{C}'\) be a second complex, \((\mathbf{F}^p\mathbf{C}')_{p\in \mathbb{Z}}\) a decreasing sequence of subcomplexes of \(\mathbf{C}'\) , \(\mathbf{E}'\) the associated spectral sequence. Let \(u:\mathbf{C}\to \mathbf{C}'\) be a morphism of complexes such that \(u(\mathbf{F}^p\mathbf{C})\subset \mathbf{F}^p\mathbf{C}'\) for all \(p\) . Show that \(u\) induces a morphism of spectral sequences \(\tilde{u}:\mathbf{E}\to \mathbf{E}'\) .
\item
  With the notations of \(d\) ), let \(v: \mathbf{C} \to \mathbf{C}'\) be a second morphism such that \(v(\mathbf{F}^p \mathbf{C}) \subset \mathbf{F}^p \mathbf{C}'\) for all \(p\) , and let \(h\) a homotopy relating \(u\) to \(v\) . Let \(k\) be an integer \(\geqslant 0\) such that \(h(\mathbf{F}^p \mathbf{C}) \subset \mathbf{F}^{p-k} \mathbf{C}'\) for all \(p\) . Show that \(\tilde{u}_r = \tilde{v}_r\) for \(r > k\) .
\end{enumerate}

\begin{enumerate}
\def\labelenumi{\arabic{enumi})}
\setcounter{enumi}{16}
\tightlist
\item
  Let C be a bicomplex (X, p. 174, exercise 11).
\end{enumerate}

\begin{enumerate}
\def\labelenumi{\alph{enumi})}
\tightlist
\item
  One defines two decreasing sequences \(({}^{\prime}\mathbf{F}^p\mathbf{C})_{p\in \mathbf{Z}}\) and \(({}^{\prime\prime}\mathbf{F}^q\mathbf{C})_{q\in \mathbf{Z}}\) of subcomplexes of the complex associated with C by setting
\end{enumerate}

\[
\left(^ {\prime} \mathrm{F} ^ {p} \mathrm{C}\right) ^ {n} = \bigoplus_ {i \geqslant p} \mathrm{C} ^ {i, n - i} \quad \left(^ {\prime \prime} \mathrm{F} ^ {q} \mathrm{C}\right) ^ {n} = \bigoplus_ {j \geqslant q} \mathrm{C} ^ {n - j, j}.
\]

Let ′E and ″E be the corresponding spectral sequences (exercise 16, a)). Show that for every pair \((p, q)\) the A-module \({}^{\prime}E_2^{pq}\) (resp. \({}^{\prime\prime}E_2^{pq}\) ) is isomorphic to \(H^{\prime\prime p}(H^{\prime q}(C))\) (resp. \(H^{\prime p}(H^{\prime\prime q}(C))\) ).

\begin{enumerate}
\def\labelenumi{\alph{enumi})}
\setcounter{enumi}{1}
\item
  Suppose there exists an integer \(n\) such that \(C^{pq} = 0\) for \(p \geqslant n\) or for \(q \leqslant n\) . Show that the spectral sequence ′E converges to H(C). If likewise there exists an integer \(m\) such that \(C^{pq} = 0\) for \(p \leqslant m\) or for \(q \geqslant m\) , the spectral sequence ″E converges to H(C).
\item
  We assume that \(C^{pq} = 0\) for p \textless{} 0 or q \textless{} 0. Show that the homomorphism \(\delta^{p}: E_{2}^{p0} \to H^{p}(C)\) (exercise 15) comes, by passing to homology, from the canonical injection of complexes ``Z\^{}\{0\}(C) → C''. Likewise the homomorphism ``E\_\{2\}\^{}\{p0\} → H\^{}\{p\}(C)'' comes from the injection ``Z\^{}\{0\}(C) → C''.
\item
  Let \(\tilde{C}\) a second bicomplex, ′E and ″E the two associated spectral sequences. Show that every morphism of bicomplexes \(u: C \to \tilde{C}\) induces morphisms of spectral sequences ′u : ′E → ′E and ″u : ″E → ″E. Show that if two morphisms \(u\) , \(v\) of \(C\) to \(\tilde{C}\) are homotopic, the associated morphisms ′u and ′v (resp. ″u and ″v) are equal (use exercise 16, e)).
\end{enumerate}

\begin{enumerate}
\def\labelenumi{\arabic{enumi})}
\setcounter{enumi}{17}
\tightlist
\item
  Let E be a spectral sequence (exercise 14), such that the A-module \(\mathbf{E}_2\) has finite length. Let \(\mathcal{C}\) be the set of isomorphism classes of A-modules of finite length; show that one has \(\chi_{\mathcal{G}}(\mathbf{E}_2) = \chi_{\mathcal{G}}(\mathbf{E}_r) = \chi_{\mathcal{G}}(\mathbf{E}_{\infty})\) for all \(r \geqslant 2\) (X, p. 41). If moreover E converges to a graded A-module G, the latter has finite length and one has \(\chi_{\mathcal{G}}(\mathbf{G}) = \chi_{\mathcal{G}}(\mathbf{E}_2)\) .
\end{enumerate}

¶ 19) One calls a simplicial scheme a pair (K, F), where K is a set and F a set of finite subsets of K, called faces of K, such that every subset of a face is a face and every point of K belongs to at least one face. A simplicial map from a simplicial scheme (K, F) to a simplicial scheme (L, G) is a map from K to L which carries every face of K to a face of L. a) Let n be an integer ≥ 0. One calls an n-simplex of the simplicial scheme K every sequence (α₀, \ldots, αₙ) of elements of K such that the αᵢ belong to a single face of K. One denotes by Sₙ(K, A) the free A-module generated by the n-simplexes for n ≥ 0, and one sets Sₙ(K, A) = 0 for n \textless{} 0. One defines an A-homomorphism \(d_{n}:\mathbf{S}_{n}(\mathbf{K},\mathbf{A})\to \mathbf{S}_{n - 1}(\mathbf{K},\mathbf{A})\) by setting \(d_{n} = 0\) for \(n\leqslant 0\) and \(d_{n}(\alpha_{0},\dots,\alpha_{n}) = \sum_{i = 0}^{n}(-1)^{i}(\alpha_{0},\dots,\theta_{i},\dots,\alpha_{n})\)

for all \(n\) -simplex \((\alpha_0, \ldots, \alpha_n)\) (the sign over a letter means that it is to be omitted). Show that \(d_{n-1} \circ d_n = 0\) for all \(n\) , so that \((S(K, A), d)\) is a complex of A-modules. One defines in the same way an A-complex \(C(K, A)\) by setting \(C^n(K, A) = \text{Hom}_Z(S_n(K, Z), A)\) , \(\delta^n = \text{Hom}(d_n, 1_A)\) for all \(n\) . For \(n \geqslant 0\) , one sometimes denotes \(H_n(K, A)\) (resp. \(H^n(K, A)\) ) the A-module \(H_n(S(K, A))\) (resp. \(H^n(C(K, A)))\) . b) Let K, L be two simplicial schemes, \(f: K \to L\) a simplicial map. Show that \(f\) induces morphisms of complexes \(S(f): S(K, A) \to S(L, A)\) and \(C(f): C(L, A) \to C(K, A)\) , whence A-homomorphisms \(f_*: H(K, A) \to H(L, A)\) and \(f^*: H'(L, A) \to H'(K, A)\) .

\begin{enumerate}
\def\labelenumi{\alph{enumi})}
\setcounter{enumi}{2}
\tightlist
\item
  Let \(g: \mathbf{K} \to \mathbf{L}\) a second simplicial map. We say that \(f\) and \(g\) are simplicially homotopic if for every face \(\sigma\) of \(\mathbf{K}\) , the set \(f(\sigma) \cup g(\sigma)\) is a face of \(\mathbf{L}\) . Show that the morphisms \(\mathbf{S}(f)\) and \(\mathbf{S}(g)\) (resp. \(\mathbf{C}(f)\) and \(\mathbf{C}(g)\) ) are then homotopic (define a homotopy \(h\) by setting
\end{enumerate}

\[
h (\alpha_ {0}, \dots , \alpha_ {n}) = \sum_ {i = 0} ^ {n} (- 1) ^ {i} (f (\alpha_ {0}), \dots , f (\alpha_ {i}), g (\alpha_ {i}), \dots , g (\alpha_ {n}))
\]

for every n-simplex \((\alpha_{0}, \ldots, \alpha_{n})\) .

\begin{enumerate}
\def\labelenumi{\alph{enumi})}
\setcounter{enumi}{3}
\tightlist
\item
  A simplicial scheme K is said to be conical if there exists a point s of K such that for every face σ of K, σ ∪ \{s\} be a face. Show that there then exist homotopisms from S(K, A) and from C(K, A) onto A. e) Let D(K, A) denote the subcomplex of S(K, A) generated by the simplices (α₀, \ldots, αₙ) for which two of the αᵢ are equal, and by the elements of the form (α₀, \ldots, αₙ) - ε(σ). (ασ(0), \ldots, ασ(n)) for every n-simplex (α₀, \ldots, αₙ) and every permutation σ of \{0, \ldots, n\}Show that D(K, A) is a subcomplex of S(K, A), and that the quotient map p : S(K, A) → S(K, A)/D(K, A) is a homotopy equivalence (a homotopy inverse is obtained by choosing a total order on K and associating to the class modulo D(K, A) of a simplex (α₀, \ldots..., αₙ) the simplex ε(σ). (ασ(0), \ldots, ασ(n)), where σ is the permutation such that ασ(0) \textless{} \ldots{} \textless{} ασ(n), if all the αᵢ are distinct, and 0 otherwise).
\end{enumerate}

*20) Let R be a complete discrete valuation ring (AC, VI, § 3, no. 6) with residue field k of characteristic p \textgreater{} 0. Assume that the field of fractions of R has characteristic 0. Let v denote the normalized valuation of R, m the maximal ideal of R, and R* the multiplicative group of units of R. For every integer \(n \geqslant 0\) Let R*(n) be the subgroup of R* consisting of the elements x such that \(v(x - 1) \geqslant n\) . a) Let r be an integer, \(r > v(p)/(p - 1)\) . Show that the series

\[
\exp (x) = \sum_ {i = 0} ^ {+ \infty} x ^ {i} / i!
\]

converges for \(x \in \mathfrak{m}^r\) and defines an isomorphism from the additive group \(\mathfrak{m}^r\) onto the group \(\mathbf{R}^*(r)\) . b) Suppose that \(k\) is a finite field with \(q\) elements. Let \(n\) be an integer \(\geqslant 1\) , \(\mu_n\) the subgroup of elements \(x\) of \(\mathbf{R}^*\) such that \(x^n = 1\) . Show that the groups \(\mu_n\) and \(\mathbf{R}^*/(\mathbf{R}^*)^n\) are finite and that one has

\[
\operatorname{Card} \left(\mathrm{R} ^ {*} / \mathrm{R} ^ {* n}\right) / \operatorname{Card} \left(\mu_ {n}\right) = q v (n).
\]

(Let \(\mathcal{C}\) the set of classes of finite \(\mathbf{Z}\) -modules, \(\varphi : \mathrm{K}(\mathcal{C}) \to \mathbf{Q}^*\) be the homomorphism defined by

\[
\varphi ([ \mathrm{G} ]) = \text { Card } (\mathrm{G}) \quad \text { for } \quad \mathrm{G} \in \mathscr {C},
\]

\(C^{(r)}\) the complex of Z-modules such that \(C_{p}^{(r)} = 0\) for \(p \neq 0, 1\) , \(C_{0}^{(r)} = C_{1}^{(r)} = R^{*}(r)\) , \(d_{1}(x) = x^{n}\) for \(x \in R^{*}(r)\) . Prove that \(\varphi(H(C^{(r)}))\) is independent of r, then use a).

\begin{enumerate}
\def\labelenumi{\arabic{enumi})}
\setcounter{enumi}{20}
\tightlist
\item
  Assume that the ring A is a commutative algebra over a commutative ring k; the A-module \(D_{k}(A)\) of k-derivations of A is then identified with the dual of the A-module \(\Omega_{A/k}^{1}\) , which defines on \(\Omega_{A/k}\) a structure of \(\Lambda_{\mathbf{A}}(\mathbf{D}_{k}(\mathbf{A}))\) left module (cf. III, p. 165). Prove the formula
\end{enumerate}

\[
\begin{array}{l} \left(\mathbf {X} _ {0} \wedge \dots \wedge \mathbf {X} _ {p}\right) \rfloor d \omega = \sum_ {i = 0} ^ {p} (- 1) ^ {i} \mathbf {X} _ {i} \left(\left(\mathbf {X} _ {0} \wedge \dots \wedge \hat {\mathbf {X}} _ {i} \wedge \dots \wedge \mathbf {X} _ {p}\right) \rfloor \omega\right) \\ \quad + \sum_ {0 \leqslant i <   j \leqslant p} (- 1) ^ {i + j} \left([ \mathbf {X} _ {i}, \mathbf {X} _ {j} ] \wedge \mathbf {X} _ {0} \wedge \dots \wedge \hat {\mathbf {X}} _ {i} \wedge \dots \wedge \hat {\mathbf {X}} _ {r} \wedge \dots \wedge \mathbf {X} _ {p}\right) \rfloor d \omega . \end{array}
\]

where \(\omega\in\Omega_{\mathbf{A}/k}^{p},\mathbf{X}_{0},\ldots,\mathbf{X}_{p}\in\mathbf{D}_{k}(\mathbf{A})\) and where the sign \(^{\wedge}\) on a letter means that it is to be omitted. (Reduce to the case \(p=1\) .)

\begin{enumerate}
\def\labelenumi{\arabic{enumi})}
\setcounter{enumi}{21}
\tightlist
\item
  We keep the conventions of the preceding exercise; if M is an A-module, \(\nabla\) a connection on M and if \(X \in D_k(A)\) , we denote by \(\nabla_X\) the \(k\) -endomorphism \((1_M \otimes X) \circ \nabla\) of M.
\end{enumerate}

\begin{enumerate}
\def\labelenumi{\alph{enumi})}
\item
  Let \(\mathbf{M}_1, \mathbf{M}_2\) two A-modules, equipped with connections \(^1\nabla\) , \(^2\nabla\) ; show that one defines a connection \(\nabla\) over \(\mathbf{M}_1 \oplus \mathbf{M}_2\) (resp. \(\mathbf{M}_1 \otimes_{\mathbf{A}} \mathbf{M}_2\) , resp. \(\operatorname{Hom}_{\mathbf{A}}(\mathbf{M}_1, \mathbf{M}_2)\) ) by setting \(\nabla(m_1 + m_2) = ^1\nabla(m_1) + ^2\nabla(m_2)\) for \(m_1 \in \mathbf{M}_1\) , \(m_2 \in \mathbf{M}_2\) (resp. \(\nabla_{\mathbf{X}}(m_1 \otimes m_2)) = ^1\nabla_{\mathbf{X}}(m_1) \otimes m_2 + m_1 \otimes ^2\nabla_{\mathbf{X}}(m_2)\) for \(\mathbf{X} \in \mathbf{D}_k(\mathbf{A})\) , \(m_1 \in \mathbf{M}_1\) and \(m_2 \in \mathbf{M}_2\) , resp. \(\nabla_{\mathbf{X}}(u) = ^2\nabla_{\mathbf{X}} \circ u - u \circ ^1\nabla_{\mathbf{X}}\) for \(\mathbf{X} \in \mathbf{D}_k(\mathbf{A})\) , \(u \in \operatorname{Hom}_{\mathbf{A}}(\mathbf{M}_1, \mathbf{M}_2)\) .
\item
  Let M be an A-module, \(\nabla: \mathbf{M} \to \mathbf{M} \otimes_{\mathbf{A}} \Omega_{\mathbf{A}/k}^{1}\) a connection on M. If X, Y ∈ D \(_{k}\) (A), we denote by R(X, Y) the A-endomorphism of M such that R(X, Y)(m) = \textless R(m), X ∧ Y\textgreater{} for all m ∈ M. Prove the formula: R(X, Y) = {[}∇X, ∇Y{]} - ∇{[}X,Y{]} (use Exercise 21).
\item
  Assume the A-module M is projective, so that the curvature homomorphism R is identified with an element of \(\operatorname{End}_{\mathbf{A}}(\mathbf{M}) \otimes_{\mathbf{A}} \Omega_{\mathbf{A}/k}^{2}\) . If one equips \(\operatorname{End}_{\mathbf{A}}(\mathbf{M})\) with the connection \(\nabla\) defined in \(a\) ), show that one has \(\nabla^{2}\mathbf{R} = 0\) .
\end{enumerate}

